\paragraph{総消費水準目標（CLT）}
\label{sec:chapter5_beta_shock_policies_clt}
式
\eqref{form:chapter3_policies_monetary_home_clt_i_H_eq}、
\eqref{form:chapter3_policies_monetary_home_clt_gamma_H_eq}および
\eqref{form:chapter3_policies_monetary_home_clt_chi_H_eq}
により記述されたとおり、CLTの式は以下のようであった。
\begin{align}
    \label{form:chapter5_beta_shock_policies_clt_i_H_eq}
        i_t^H&=i_{ss}^H+\phi_{gap}^H(\log(c_t^{H\to W})-\log(\chi_t^H))+\phi_{level}^H\gamma_t^H \\
    \label{form:chapter5_beta_shock_policies_clt_gamma_H_eq}
        \gamma_t^H&=\gamma_{t-1}^H+(\log(c_{t-1}^{H\to W})-\log(\chi_{t-1}^H)) \\
    \label{form:chapter5_beta_shock_policies_clt_chi_H_eq}
        \log(\chi_t^H)-\log(\chi_{ss}^H)&=\rho_{\chi}^H\left(\log(\chi_{t-1}^H)-\log(\chi_{ss}^H)\right)+\varepsilon_t^{\chi,H}
\end{align}
ここで
\(\phi_{gap}^H,\phi_{level}^H\)は反応係数、
\(\chi_t^H\)は目標パス、
\(\gamma_t^H\)は過去の乖離の累積、
\(\varepsilon_t^{\chi,H}\)は外生ショック項である。

\subparagraph{\(i_t^H\)}
\label{sec:chapter5_beta_shock_policies_clt_i_H}     
名目利子率\(i_t^H\)の図
\ref{fig:chapter5_beta_shock_graphs_i_H}
を見ると、CLTはショック発生直後の\(t=0\)において名目利子率を即座に下限まで引き下げていることがわかる。
他の多くの政策が時間をかけて徐々に利子率を下げているのに対し、
CLTは初動の引き下げ幅が最も大きく、序盤から中盤にかけて下限に張り付くような動きを見せている。
その後は120期付近で下限を脱し、プラス圏に浮上することなく緩やかに定常状態へと回帰する軌道を描いている。

\subparagraph{\(\operatorname{E}_t[\lambda_{t+1}^H]\)}
\label{sec:chapter5_beta_shock_policies_clt_E_lambda_H}
所得の限界効用\(\lambda_t^H\)の図
\ref{fig:chapter5_beta_shock_graphs_lambda_H}
を見ると、CLTはショック発生直後の\(t=0\)においてプラス圏からスタートし、
その後は急速に低下して50期付近でマイナス圏の底を打つ軌道を描いている。
他の多くの政策が序盤の低下の後に定常状態であるゼロ付近へと収束、あるいはマイナス圏で推移するのに対し、
CLTは100期付近でゼロを上抜けた後もプラス圏を緩やかに上昇し続け、
定常状態から持続的に遠ざかる特異な動きを見せている。

\subparagraph{\(\operatorname{E}_t[p_{t+1}^H]\)}
\label{sec:chapter5_beta_shock_policies_clt_E_p_H}
生産者物価指数\(p_t^H\)の図
\ref{fig:chapter5_beta_shock_graphs_p_H}
を見ると、CLTはショック発生直後から一貫してマイナス圏を右肩下がりで推移していることがわかる。
他の政策が定常状態であるゼロ近傍で安定するか、あるいはプラス圏へと向かっていくのに対し、
CLTは期間を通じて際限なく下落し続けている。
400期に至ってもグラフの傾きは平坦にならず、
定常状態へ回帰する兆しを全く見せないまま下方に大きく乖離していく特異な軌道を描いている。

\subparagraph{\(c_t^{H\to W}\)}
\label{sec:chapter5_beta_shock_policies_clt_c_H_W}
総消費指数\(c_t^{H\to W}\)の図
\ref{fig:chapter5_beta_shock_graphs_c_H_W}
を見ると、CLTは前半において他のどの政策よりも落ち込みを小さく抑え込んでいることがわかる。
ショック直後の不況期において消費水準を最も高く維持できているため、
この前半部分においては厚生を強力に改善する効果を発揮している。
一方で後半に入ると、CLTは目標である初期定常状態の水準（ゼロ）へと速やかに回帰する。
そのため、消費者物価水準目標（CPLT）や消費者物価インフレ目標（IT-CPI）のような
後半の巨大なオーバーシュート（プラス方向への乖離）を持たず、
中盤以降の消費増加による厚生の押し上げ効果は他の政策よりも小さくなっている。

\subparagraph{\(y_t^H\)}
\label{sec:chapter5_beta_shock_policies_clt_y_H}
生産\(y_t^H\)の図
\ref{fig:chapter5_beta_shock_graphs_y_H}
を見ると、CLTの生産は総消費指数\(c_t^{H\to W}\)と同様に
「前半には他の政策よりも多いが、後半には他の政策よりも少ない」という波形を描いていることがわかる。
ショック直後の前半において、CLTは生産の落ち込みを他のどの政策よりも小さく防いでいる。
生産の維持は労働の高止まりを意味するため、この前半部分においては
労働の非効用を通じて厚生を悪化させる効果が他の政策よりも大きくなっている。
一方で中盤以降のプラス圏での推移（オーバーシュート）を見ると、
CLTは定常状態へ速やかに回帰するため他の政策よりも低い水準で推移する。
消費者物価水準目標（CPLT）や消費者物価インフレ目標（IT-CPI）のような巨大な山を作らないため、
後半における生産過剰に伴う厚生のマイナス効果を最小限に防いでいる。

\subparagraph{\(\pi_t^H\)}
\label{sec:chapter5_beta_shock_policies_clt_pi_H}
グロス・インフレ率\(\pi_t^H\)の図
\ref{fig:chapter5_beta_shock_graphs_pi_H}を見ると、
CLTの定常状態からの乖離は全政策の中で最も大きく、極めて不安定な軌道を描いていることがわかる。
ショック直後の前半には他のどの政策よりも大きくプラス（インフレ）へ振れ、
中盤以降は一転して大きくマイナス（デフレ）へと落ち込み続けている。
名目アンカーを持たず実質変数のみを目標としているためインフレ率の変動が全く制御されておらず、
この継続的かつ極端なインフレ率の乖離が厚生を大きく悪化させる最大の要因となっている。

\subparagraph{\(\Delta_{t-1}^H\)}
\label{itm:chapter5_beta_shock_policies_clt_Delta_t_minus_1_H}
価格分散\(\Delta_{t-1}^H\)の図
\ref{fig:chapter5_beta_shock_welfare_Delta_H}を見ると、
CLTの定常状態からの乖離は全政策の中で最も著しく、期間を通じて増大し続けていることがわかる。
他の政策がゼロ近傍を維持するか、あるいは一時的な山を形成した後に収束へと向かうのに対し、
CLTは名目アンカーを持たないためにインフレ率の乖離が永続的に蓄積し、
価格分散が際限なく拡大していく軌道を描いている。
この価格分散の圧倒的な増大による生産効率の悪化が、CLTの厚生を大きく押し下げる致命的な要因となっている。